% ===============================================
% MATH 3053: Abstract Algebra I         Fall 2017
% hw_revised_cn.tex
% 中文适配版作业修订提交模板
% ===============================================
%         请仔细阅读以下说明！！！
% ===============================================
% 生成PDF提交时，将文件名改为 hwX-姓名.pdf
% X替换为作业编号，姓名替换为你的姓氏
% ===============================================


% -----------------------------------------------
% 以下序言部分可忽略，直接跳至 "START HERE" 区域
% -----------------------------------------------

\documentclass{article}
% 页面设置与基础宏包
\usepackage[margin=1in]{geometry}  % 页边距1英寸
\usepackage{amsmath,amsthm,amssymb} % 数学公式、定理环境、符号
\usepackage{ctex}   % 中文支持核心宏包（Overleaf预装）
\usepackage{pifont}                 
\renewcommand{\qedsymbol}{}
% 常用数学集合符号定义
\newcommand{\R}{\mathbb{R}}  % 实数集
\newcommand{\Z}{\mathbb{Z}}  % 整数集
\newcommand{\N}{\mathbb{N}}  % 自然数集
\newcommand{\Q}{\mathbb{Q}}  % 有理数集
\newcommand{\C}{\mathbb{C}}  % 复数集

% 定理、习题等环境定义（支持中英文标注）
\newenvironment{theorem}[2][Theorem]{\begin{trivlist}
\item[\hskip \labelsep {\bfseries #1}\hskip \labelsep {\bfseries #2.}]}{\end{trivlist}}
\newenvironment{lemma}[2][Lemma]{\begin{trivlist}
\item[\hskip \labelsep {\bfseries #1}\hskip \labelsep {\bfseries #2.}]}{\end{trivlist}}
\newenvironment{exercise}[2][Exercise]{\begin{trivlist}
\item[\hskip \labelsep {\bfseries #1}\hskip \labelsep {\bfseries #2.}]}{\end{trivlist}}
% 中文"问题"环境（与原problem并存，可按需选用）
\newenvironment{problem}[2][Problem]{\begin{trivlist}
\item[\hskip \labelsep {\bfseries #1}\hskip \labelsep {\bfseries #2.}]}{\end{trivlist}}
\newenvironment{problem_cn}[2][问题]{\begin{trivlist}
\item[\hskip \labelsep {\bfseries #1}\hskip \labelsep {\bfseries #2.}]}{\end{trivlist}}
\newenvironment{question}[2][Question]{\begin{trivlist}
\item[\hskip \labelsep {\bfseries #1}\hskip \labelsep {\bfseries #2.}]}{\end{trivlist}}
\newenvironment{corollary}[2][Corollary]{\begin{trivlist}
\item[\hskip \labelsep {\bfseries #1}\hskip \labelsep {\bfseries #2.}]}{\end{trivlist}}

% 解答环境（支持中英文标注）
\newenvironment{solution}{\begin{proof}[Solution]}{\end{proof}}
\newenvironment{solution_cn}{%
  \begin{proof}[解答]%
  \kaishu  % 切换中文为楷体（ctex 提供的命令）
}{%
  \end{proof}%
}

\begin{document}

% ------------------------------------------ %
%                 从这里开始编辑                 %
% ------------------------------------------ %

% 标题与作者信息（中文适配）
\title{《优化方法基础》课程作业\\无约束凸优化问题}  % X替换为作业编号
\author{姓名：林圣翔 \quad 班级：计试2201 \quad 学号：2223312202}  % 替换括号内容

\maketitle

\begin{problem_cn}{1}
假设函数$f$是满足$mI \leq \nabla^2 f(x) \leq MI$的强凸函数。令$d$为$x$处的下降方向。证明对于  
$$
0 < t \leq -\frac{\nabla f(x)^\top d}{M\|d\|_2^2}
$$  
回溯终止条件能够满足。
\end{problem_cn}
\begin{solution_cn} 
\quad \\
$\because$ \( f \) 是强凸函数且 Hessian 矩阵有上界 \( M \) \\
$\therefore$ 根据Taylor展开和二次上界性质，有  
$$
f({x}^k + t{d}^k) \leq f({x}^k) + t\nabla f({x}^k)^\top {d}^k + \frac{M}{2} t^2 \|{d}^k\|_2^2
$$  
$\because$ 要满足回溯终止条件，只需证明
$$
f({x}^k) + t\nabla f({x}^k)^\top {d}^k + \frac{M}{2} t^2 \|{d}^k\|_2^2 \leq f({x}^k) + \alpha t\nabla f({x}^k)^\top {d}^k
$$
即
$$
t\nabla f({x}^k)^\top {d}^k + \frac{M}{2} t^2 \|{d}^k\|_2^2 \leq \alpha t\nabla f({x}^k)^\top {d}^k
$$
$\because$ \( t > 0 \) 和 \( \nabla f({x}^k)^\top {d}^k < 0 \)，两边除以 \( t \)：  
$$
\nabla f({x}^k)^\top {d}^k + \frac{M}{2} t \|{d}^k\|_2^2 \leq \alpha \nabla f({x}^k)^\top {d}^k
$$  
即 
$$
\frac{M}{2} t \|{d}^k\|_2^2 \leq (\alpha - 1)\nabla f({x}^k)^\top {d}^k
$$  
$\because$ \( \alpha -1< 0 \) \\
$\therefore$
$$
t \leq \frac{2(\alpha -1)\nabla f({x}^k)^\top {d}^k}{M \|{d}^k\|_2^2}
$$  
$\because$ 取\( \alpha = 0.5 \)，代入可得
$$
t \leq -\frac{\nabla f({x}^k)^\top {d}^k}{M \|{d}^k\|_2^2}
$$  
$\therefore$当 \( 0 < t \leq -\dfrac{\nabla f({x}^k)^\top {d}^k}{M \|{d}^k\|_2^2} \) 时，回溯终止条件成立。  
\end{solution_cn}
\vspace{0.25in}
\begin{problem_cn}{2}
min f(x)迭代求解最优性条件 $\nabla f(x^*) = 0$

\textbf{假设:}

$\checkmark$ 存在 $m,M$ 满足 $mI \leq \nabla^2 f(x) \leq MI$ $(\forall x \in S)$  
    
$\checkmark$ 存在常数 $\gamma, \tilde{\gamma} \in (0,1]$，使得 $\|x\| \geq \gamma \|x\|_2$，$\|x\|_* \geq \tilde{\gamma} \|x\|_2$  

\textbf{证明：}

1. 精确直线搜索时，非规范化最速下降方法的收敛性；  

2. 回溯直线搜索时，非规范化最速下降方法的收敛性；  
\end{problem_cn}
\begin{solution_cn}
\quad \\
1.证明.精确直线搜索时，非规范化最速下降方法的收敛性\\
在精确直线搜索中，步长 \( t \) 被选择为最小化 \( f({x}^k + t{d}^k) \)。利用 \( f \) 的光滑性，有：  
$$
f({x}^k + t{d}^k) \leq f({x}^k) + t\nabla f({x}^k)^\top {d}^k + \frac{M}{2} t\|{d}^k\|_2^2.
$$
$\because$\( \|{d}^k\| \leq 1 \) , \( \|{x}\| \geq \gamma\|{x}\|_2 \) \\
$\therefore$ \( \|{d}^k\|_2 \leq \frac{1}{\gamma} \)
代入 \( \nabla f({x}^k)^\top {d}^k = -\|\nabla f({x}^k)\|_* \) 得
$$
f({x}^k + t{d}^k) \leq f({x}^k) - t\|\nabla f({x}^k)\|_* + \frac{M}{2\gamma^2} t^2.
$$  
$\because$ 通过精确直线搜索，最小化右侧关于 \( t \) 的二次函数，得到最优步长 \( t^* = \frac{\gamma^2 \|\nabla f({x}^k)\|_*}{M} \)\\
$\therefore$ 可得
$$
f({x}^{k+1}) \leq f({x}^k) - \frac{\gamma^2}{2M} \|\nabla f({x}^k)\|_*^2.
$$  
$\because$ \( \|\nabla f({x}^k)\|_* \geq \tilde{\gamma} \|\nabla f({x}^k)\|_2 \) ,\(\|\nabla f({x}^k)\|_2^2 \geq 2m(f({x}^k) - p^*) \)\\
$\therefore$ 可得
$$
f({x}^{k+1}) - p^* \leq f({x}^k) - p^* - \frac{\gamma^2}{2M} \tilde{\gamma}^2 \cdot 2m(f({x}^k) - p^*) = \left(1 - \frac{\gamma^2 \tilde{\gamma}^2 m}{M}\right)(f({x}^k) - p^*)
$$  
令 \( c = 1 - \frac{\gamma^2 \tilde{\gamma}^2 m}{M} < 1 \)，则
$$
f({x}^k) - p^* \leq c^k (f({x}^0) - p^*)
$$
$\therefore$精确直线搜索时，非规范化最速下降方法具备收敛性\\
2.证明：回溯直线搜索时，非规范化最速下降方法的收敛性\\
在回溯直线搜索中，步长 \( t \) 从 \( t = 1 \) 开始，并以因子 \( \beta \in (0,1) \) 减少，直到满足回溯终止条件
$$
f({x}^k + t{d}^k) \leq f({x}^k) + \alpha t \, \nabla f({x}^k)^\top {d}^k = f({x}^k) - \alpha t \, \|\nabla f({x}^k)\|_*,
$$  
其中 \( \alpha \in (0,1) \)。利用光滑性 
$$
f({x}^k + t{d}^k) \leq f({x}^k) - t\|\nabla f({x}^k)\|_* + \frac{M}{2\gamma^2} \, t^2
$$
$\because$为确保回溯终止条件满足，有
$$
f({x}^k) - t\|\nabla f({x}^k)\|_* + \frac{M}{2\gamma^2} \, t^2 \leq f({x}^k) - \alpha t \, \|\nabla f({x}^k)\|_*
$$
$\therefore$简化得
$$
\frac{M}{2\gamma^2} \, t^2 \leq (1 - \alpha) t \, \|\nabla f({x}^k)\|_*
$$  
$\therefore$当 \( t \leq \dfrac{2\gamma^2 (1 - \alpha) \|\nabla f({x}^k)\|_*}{M} \) 时，回溯终止条件成立。在回溯中，接受的 \( t \) 满足：  
$$
t \geq \dfrac{2\gamma^2 (1 - \alpha) \|\nabla f({x}^k)\|_*}{M}
$$
$\therefore$
$
f({x}^{k+1}) \leq f({x}^k) - \alpha t \|\nabla f({x}^k)\|_* \leq f({x}^k) - \frac{2\alpha\gamma^2(1 - \alpha)}{M} \|\nabla f({x}^k)\|_2
$ \\
$\because$ \( \|\nabla f({x}^k)\|_* \geq \tilde{\gamma} \|\nabla f({x}^k)\|_2 \) ,\(\|\nabla f({x}^k)\|_2^2 \geq 2m(f({x}^k) - p^*) \)\\
$\therefore$ 可得
$$
f({x}^{k+1}) - p^* \leq f({x}^k) - p^* - \frac{2\alpha\gamma^2(1 - \alpha)}{M} \tilde{\gamma}^2 \cdot 2m(f({x}^k) - p^*) = \left(1 - \frac{2\alpha\gamma^2(1 - \alpha)\tilde{\gamma}^2 m}{M}\right)(f({x}^k) - p^*)
$$  
令 \( c = 1 - \dfrac{2\alpha\gamma^2(1 - \alpha)\tilde{\gamma}^2 m}{M} < 1 \)，则  
$$
f({x}^k) - p^* \leq c^k (f({x}^0) - p^*)
$$  
$\therefore$回溯直线搜索时，非规范化最速下降方法具备收敛性 
\end{solution_cn}
\vspace{0.25in}
\begin{problem_cn}{3}
min f(x)迭代求解最优性条件 $\nabla f(x^*) = 0$

\textbf{假设:}

$\checkmark$ 存在 $m, M$ 满足 $mI \leq \nabla^2 f(x) \leq MI$ $\left( \forall x \in S \right)$  

$\checkmark$ 存在 $L$ 满足 $\big\|\nabla^2 f(y) - \nabla^2 f(x)\big\|_2 \leq L\|y - x\|_2$  

\textbf{证明：} 回溯直线搜索时牛顿方法的收敛性。
\end{problem_cn}
\begin{solution_cn}
\quad \\
1.阻尼牛顿阶段\\
当 $\|\nabla f({x}^k)\|_2 \geq \eta$ 时（其中 $\eta > 0$ 是一个常数），函数值以线性速率下降。\\  
由泰勒展开可得
$$
f({x}^k + t {d}_k^n) \leq f({x}^k) + t \nabla f({x}^k)^\top {d}_k^n + \frac{1}{2} M t^2 \|{d}_k^n\|_2^2
$$
代入$\nabla f({x}^k)^\top {d}_k^n = -\lambda({x}^k)^2$ 和 $\|{d}_k^n\|_2^2 \leq \frac{1}{m} \lambda({x}^k)^2$（因 $\|{d}_k^n\|_2 \leq \frac{1}{\sqrt{m}} \lambda({x}^k)$），得
$$
f({x}^k + t {d}_k^n) \leq f({x}^k) - t \lambda({x}^k)^2 + \frac{M}{2m} t^2 \lambda({x}^k)^2
$$  
$\because$回溯终止条件要求
$$
f({x}^k + t {d}_k^n) \leq f({x}^k) - \alpha t \lambda({x}^k)^2
$$  
$\therefore$需满足
$$
-t \lambda({x}^k)^2 + \frac{M}{2m} t^2 \lambda({x}^k)^2 \leq -\alpha t \lambda({x}^k)^2
$$  
简化得  
$$
\frac{M}{2m} t \leq 1 - \alpha \implies t \leq \frac{2m(1 - \alpha)}{M}
$$
$\because$在回溯搜索中，$t$ 被选为满足 $t \geq \beta \cdot \frac{2m(1 - \alpha)}{M}$（$\beta \in (0,1)$ 为回溯因子）。\\
$\therefore$
$$
f({x}^{k+1}) - f({x}^k) \leq -\alpha t \lambda({x}^k)^2 \leq -\alpha \beta \frac{2m(1 - \alpha)}{M} \lambda({x}^k)^2
$$  
$\because$ $\lambda({x}^k)^2 \geq \frac{1}{M} \|\nabla f({x}^k)\|_2^2$ 且 $\|\nabla f({x}^k)\|_2 \geq \eta$\\
$\therefore$
$$
f({x}^{k+1}) - f({x}^k) \leq -\alpha \beta \frac{2m(1 - \alpha)}{M^2} \eta^2 = -\gamma
$$
其中 $\gamma = \alpha \beta \frac{2m(1 - \alpha)}{M^2} \eta^2 > 0$。\\ 
$\because$ $f({x}) \geq p^*$（$p^*$ 为最优值）\\
$\therefore$ 阻尼牛顿阶段的迭代次数最多为 $\dfrac{f({x}^0) - p^*}{\gamma}$。\\
2.二次收敛阶段\\
当 $\|\nabla f({x}^k)\|_2 < \eta$ 时（其中 $\eta$ 足够小），回溯搜索接受 $t=1$，且收敛速度为二次。\\
由泰勒展开可得
$$
f({x}^k + {d}_k^n) \leq f({x}^k) + \nabla f({x}^k)^\top {d}_k^n + \frac{1}{2} ({d}_k^n)^\top \nabla^2 f({x}^k) {d}_k^n + \frac{L}{6} \|{d}_k^n\|_2^3
$$  
代入以下关系：  
$$\nabla f({x}^k)^\top {d}_k^n = -\lambda({x}^k)^2
$$  
$$({d}_k^n)^\top \nabla^2 f({x}^k) {d}_k^n = \lambda({x}^k)^2
$$  
$$\|{d}_k^n\|_2^3 \leq \dfrac{1}{m^{3/2}} \lambda({x}^k)^3
$$
可得
$$
f({x}^k +{d}_k^n) \leq f({x}^k) - \lambda({x}^k)^2 + \frac{1}{2} \lambda({x}^k)^2 + \frac{L}{6m^{3/2}} \lambda({x}^k)^3 = f({x}^k) - \frac{1}{2} \lambda({x}^k)^2 + \frac{L}{6m^{3/2}} \lambda({x}^k)^3
$$
$\because$回溯终止条件要求
$$
f({x}^k + {d}_k^n) \leq f({x}^k) - \alpha \lambda({x}^k)^2 \quad (\alpha \in (0,1))
$$  
$\therefore$需满足不等式
$$
-\frac{1}{2} \lambda({x}^k)^2 + \frac{L}{6m^{3/2}} \lambda({x}^k)^3 \leq -\alpha \lambda({x}^k)^2
\implies
\frac{L}{6m^{3/2}} \lambda({x}^k) \leq \frac{1}{2} - \alpha \implies \lambda({x}^k) \leq \frac{6m^{3/2} \left( \frac{1}{2} - \alpha \right)}{L}
$$  
利用 $\lambda({x}^k) \leq \dfrac{1}{\sqrt{m}} \|\nabla f({x}^k)\|_2$，代入上式得
$$
\|\nabla f({x}^k)\|_2 \leq \frac{6m^2 \left( \frac{1}{2} - \alpha \right)}{L}
$$  
取 $\eta = \dfrac{6m^2 \left( \frac{1}{2} - \alpha \right)}{L}$，则当 $\|\nabla f({x}^k)\|_2 < \eta$ 时，$t=1$ 满足回溯终止条件。\\
现在证明二次收敛\\
考虑
$$
\nabla f({x}^{k+1}) = \nabla f\left({x}^k + {d}_k^n\right) = \int_0^1 \left[ \nabla^2 f({x}^k + s{d}_k^n) - \nabla^2 f({x}^k) \right] {d}_k^n \, ds
$$  
$\because$ $\|\nabla^2 f({y}) - \nabla^2 f({x})\| \leq L\|{y}-{x}\|$\\
$\therefore$ $\left\|\left[\nabla^2 f({x}^k + s{d}_k^n) - \nabla^2 f({x}^k)\right] {d}_k^n\right\|_2 \leq L s \|{d}_k^n\|_2^2$ \\
$\therefore$  
$$
\|\nabla f({x}^{k+1})\|_2 \leq \int_0^1 L s \|{d}_k^n\|_2^2 \, ds = \frac{L}{2} \|{d}_k^n\|_2^2
$$  
$\because$ $\|{d}_k^n\|_2^2 \leq \frac{1}{m} \|\nabla f({x}^k)\|_2$\\
$\therefore$ 代入得
$$ 
\|\nabla f({x}^{k+1})\|_2 \leq \frac{L}{2m^2} \|\nabla f({x}^k)\|_2^2
$$  
$\therefore$ 一旦 $\|\nabla f({x}^k)\|_2 < \eta$，序列 $\{\|\nabla f({x}^k)\|_2\}$ 以二次速度收敛。\\
$\because$$f$ 是强凸函数，满足 $f({x}) - p^* \leq \frac{1}{2m} \|\nabla f({x})\|_2^2$\\
$\therefore$ $f({x}^k) - p^*$ 也快速收敛。\\
$\therefore$ 达到 $f({x}^k) - p^* < \varepsilon$ 所需的迭代次数为 $\log_2(\log_2(1/\varepsilon))$ 数量级。\\  
综上，牛顿方法在回溯直线搜索下全局收敛，且最终收敛速度为二次。  
\end{solution_cn}

% -----------------------------------------------
% 以下内容无需编辑
% -----------------------------------------------

\end{document}